\documentclass[report.tex]{subfiles}

\begin{document}

\section{Related Work}

% Proposed Length: 0.25 pages
\subsection{Verified Compilers \& Simulation Frameworks}
% Comparison with CompCert (Leroy) and CakeML (Tan et al.).
% Contrasting traditional explicit-measure stuttering with coinductive step-indexed simulation in Rocq.

% This leads to the \emph{safe backward simulation} property:
% \begin{equation*}
%   \Safe{s} \implies \forall b \in t, \, b \in s
% \end{equation*}
% \emph{Safe backward simulation} ensures specification preservation
% whenever the source program is safe,
% while granting the compiler freedom
% to remove undefined operations in unsafe programs.

% In contrast to backward simulation,
% a forward simulation asserts that every source behavior
% is reproduced by the target program: $\forall b, b \in s \implies b \in t$.
% Forward simulation can similarly be restricted to safe executions.
% Forward simulations are often easier to prove in practice,
% as reasoning proceeds along the execution steps of the source program.

% When both the source and target languages are deterministic,
% forward simulation and backward simulation coincide.
% This is how CompCert guarantees the specification preservation property.
% It exploits target determinism to derive backward simulation
% from a forward simulation proof.
% However,
% because our framework aims to accommodate non-deterministic execution,
% we directly target a safe backward simulation property.


% Proposed Length: 0.25 pages
\subsection{Relational Logics \& Program Transformation Verification}
% Comparison with Simuliris (Gäher et al.), Tail Modulo Cons (Allain et al.), and Interaction Trees (Xia et al.).
% Bringing relational separation logic directly to low-level assembly-like languages (RTL).
\end{document}

% Local Variables:
% TeX-engine: luatex
% End:

% LocalWords:  Coinductive coinductive asymmetricity arities
